\section*{\centering Resumo}

Este trabalho investiga a aplicação de \textit{retrieval-augmented generation} (RAG) no desenvolvimento de um \textit{chatbot} legislativo para consulta a discursos da 56\textordfeminine{} Legislatura do Senado Federal (2019--2023). Parte-se da dificuldade de recuperar informações parlamentares por meio de linguagem natural com precisão factual, contextualização e rastreabilidade. A metodologia consistiu na construção de um \textit{pipeline} com preparação de dados, segmentação textual, enriquecimento por metadados, geração de \textit{embeddings}, indexação vetorial e fluxo de consulta baseado em LLM. A solução foi implementada em ambiente reprodutível, com \textit{scripts} próprios para ingestão e avaliação, e validada por três estudos de caso manuais, três rodadas automatizadas com rubrica e LLM como juiz, além de revisão humana amostral das respostas automatizadas. Os resultados indicam bom desempenho em perguntas factuais, temáticas e de autoria bem delimitada, com médias entre 9,15 e 9,55 nas rodadas automatizadas. Observou-se maior instabilidade em consultas que exigem precisão numérica, comparação entre autores, controle de escopo ou síntese multifocal. Conclui-se que a abordagem é viável para recuperação auditável de discursos legislativos, desde que acompanhada de governança de fontes, avaliação sistemática e supervisão humana.

\textbf{Palavras-chave:} RAG; inteligência artificial; \textit{chatbot} legislativo; recuperação de informação; Senado Federal; LLM.
